\documentclass[11pt,letterpaper]{article}
\usepackage[utf8]{inputenc}
\usepackage[T1]{fontenc}
\usepackage[spanish,es-tabla]{babel}
\usepackage{lmodern}
\usepackage{geometry}
\usepackage{amsmath,amssymb}
\usepackage{booktabs}
\usepackage{array}
\usepackage{xcolor}
\usepackage{hyperref}
\usepackage{enumitem}
\usepackage{needspace}
\geometry{margin=2.4cm}
\hypersetup{colorlinks=true,linkcolor=blue!45!black,urlcolor=blue!55!black,citecolor=blue!45!black}
\setlist{nosep,leftmargin=*}

\title{Generación de escenarios por ajuste de momentos\\
y decisiones de portafolio en acciones chilenas\\[0.5em]
\large Protocolo e informe de investigación independiente}
\author{Gustavo Chuquiviguel}
\date{Versión 0.4 --- agosto de 2026}

\begin{document}
\maketitle

\begin{abstract}
Este informe define una investigación cuantitativa independiente sobre generación
de escenarios discretos para retornos accionarios chilenos. Se estudia una extensión
propia basada en ajuste de cuatro momentos marginales y covarianzas mediante descenso
por bloques, y se compara con distribuciones Gaussian, Student-$t$ e histórica
ponderada construidas al mismo horizonte. El criterio de éxito no es el ajuste
in-sample: son los scores probabilísticos fuera de muestra, la calibración de riesgo
y el desempeño de portafolios neto de costos frente a reglas ingenuas y robustas.
El periodo 2024--2026 ya fue observado durante el desarrollo y se declara validación,
no test confirmatorio sellado.
\end{abstract}

\section{Pregunta y contribución}

La pregunta central es si una distribución MM discreta mejora simultáneamente:
(i) el pronóstico de la distribución de retornos a cinco días; (ii) la cobertura de
medidas de cola; y (iii) las decisiones de asignación, frente a controles que reciben
exactamente la misma información.

La contribución de software es un pipeline reproducible que separa ingestión,
estimación, generación, evaluación y decisión. La contribución metodológica es
tratar el algoritmo BCD como una extensión propia falsable. No se afirma que sea una
implementación literal del algoritmo algebraico de Ponomareva, Roman y Date
\cite{ponomareva2015}, ni de la corrección de Contreras, Bosch y Herrera
\cite{contreras2018}.

\section{Antecedentes}

El ajuste de propiedades distribucionales es una familia establecida de generación de
escenarios \cite{hoyland2001,mehrotra2013}. Ponomareva et al. proponen una construcción
finita que ajusta media, covarianza y promedios marginales de momentos superiores
\cite{ponomareva2015}. Contreras et al. demuestran que una selección descuidada de
parámetros puede producir probabilidades negativas y derivan condiciones correctivas
\cite{contreras2018}. Ese resultado es especialmente pertinente porque su ejemplo usa
acciones chilenas.

Para evaluar distribuciones se necesitan scoring rules propios: un modelo no debe ser
premiado por pronósticos mal calibrados pero agudos. CRPS y Energy Score se justifican
en el marco de Gneiting y Raftery \cite{gneiting2007}. La cobertura VaR se contrasta con
Kupiec \cite{kupiec1995} y la independencia de excepciones con Christoffersen
\cite{christoffersen1998}.

\section{Datos y diseño temporal}

El universo vigente contiene 15 acciones chilenas con datos ajustados de Yahoo
Finance. El proveedor es adecuado para investigación exploratoria, no sustituye una
fuente institucional. Se conserva el panel raw, la máscara de observación y la
cobertura previa a cualquier imputación. El \textit{forward fill} está limitado a dos
días y los retornos se calculan sin relleno implícito.

La descarga vigente contiene 1.605 fechas entre 2020-01-02 y 2026-06-10, cobertura
raw del 100\% para los 15 activos y cero imputaciones. Aun así, precios sin cambio
advierten posible iliquidez: AGUAS-A, ANDINA-B, ITAUCL y PARAUCO superan una tasa
global de retornos cero del 3\%; AGUAS-A alcanza seis sesiones consecutivas en OOS.
Estas alertas no producen exclusión automática.
Los retornos cero se usan solo como proxy imperfecto de fricción, siguiendo la
motivación de Lesmond, Ogden y Trzcinka \cite{lesmond1999}; no reemplazan volumen,
spread bid--ask ni una fuente institucional de transacciones.

La muestra inicial termina el 31 de diciembre de 2023. Las fechas posteriores se usan
como validación de desarrollo. Las ventanas rolling $H=5$ sirven para estimar targets;
los scores usan ventanas no solapadas para no fingir independencia. Un test futuro
solo podrá empezar después de congelar código, hipótesis e hiperparámetros.

Como robustez temporal se ejecuta además una evaluación expanding rolling-origin en
2023, 2024, 2025 y 2026-H1. Cada origen recalibra MM y todos los controles usando solo
su historia disponible, siguiendo las recomendaciones de evaluación OOS de Tashman
\cite{tashman2000} y las salvaguardas contra leakage discutidas por Hewamalage et al.
\cite{hewamalage2023}.

\section{Modelo MM-BCD}

Sea $X\in\mathbb{R}^{N\times n}$ el soporte discreto y
$p\in\Delta_N$ su vector de probabilidades. Para momentos centrales objetivo
$M_{ki}$ y covarianza objetivo $\Sigma$, se define
\begin{equation}
 F(X,p)=\sum_{k=1}^{4}\sum_{i=1}^{n}W_{ki}
 \left(\widehat m_{ki}-M_{ki}\right)^2
 +w_\Sigma\sum_{i\ne j}A_{ij}
 \left(\widehat\Sigma_{ij}-\Sigma_{ij}\right)^2.
\end{equation}
La diagonal se excluye del segundo término porque la varianza ya se ajusta en
$k=2$. Las escalas $W_{ki}$ se construyen con potencias de volatilidad, lo que evita
divisiones inestables por medias o terceros momentos cercanos a cero.

La función realmente optimizada es
\begin{equation}
 G(X,p)=F(X,p)+\lambda_{KL}D_{KL}(p\Vert u),
 \qquad u_j=1/N.
\end{equation}
El término KL controla soluciones degeneradas sin ocultar el costo del ajuste. La
selección multi-start, el descenso y la convergencia se evalúan sobre $G$. El bloque
completo de posiciones se optimiza conjuntamente con L-BFGS porque las covarianzas
acoplan los activos; un Jacobi paralelo no es matemáticamente separable.

Los gradientes analíticos de posiciones y probabilidades están validados con
diferencias finitas direccionales. Cada subpaso se acepta solo si no aumenta $G$.

\section{Estimación temporal y dependencia}

La vida media se parametriza en semanas. Con once semanas y cinco filas rolling por
semana,
\begin{equation}
 \lambda=0.5^{1/(11\times5)}\approx0.98748,
\end{equation}
de modo que una observación de 55 filas recibe la mitad del peso de la más reciente.
El valor legacy $0.94$ tenía vida media de solo 11.2 filas.

Se aplica una contracción lineal explícita del 10\% hacia la diagonal para mejorar
condicionamiento. No se la denomina Ledoit--Wolf exacta; la literatura de shrinkage
\cite{ledoitwolf2004} motiva el control, pero la intensidad debe seleccionarse sin
usar el test final.

\Needspace{10\baselineskip}
\section{\texorpdfstring{\enspace}{}Controles y evaluación}

Los tres controles comparten media, covarianza, horizonte y ponderación:
\begin{enumerate}
 \item normal multivariada terminal;
 \item Student-$t$ multivariada terminal con covarianza parametrizada;
 \item distribución empírica EWMA.
\end{enumerate}

Las métricas primarias son CRPS marginal y Energy Score multivariado. Se reportan
además Variogram Score, error de correlación, VaR, Expected Shortfall y pruebas de
cobertura. Menor score es mejor. El ajuste de momentos dentro de muestra se conserva
como diagnóstico, no como prueba de superioridad.

La comparación sigue el principio de pérdidas predictivas pareadas de Diebold y
Mariano \cite{diebold1995}. Cada score se guarda para las mismas 121 ventanas OOS.
La dependencia temporal se conserva con moving-block bootstrap de cuatro ventanas
\cite{kunsch1989}, 5.000 remuestreos e intervalo básico centrado. La familia de nueve
contrastes por universo se ajusta con el procedimiento de Holm \cite{holm1979}.
En la evidencia pooled de varios orígenes, los bloques se remuestrean dentro de cada
fold y nunca atraviesan sus fronteras.

\section{\texorpdfstring{\enspace}{}Decisiones de portafolio}

Se comparan Equal Weight, inverse variance, HRP \cite{lopez2016}, mínima varianza
regularizada, mínimo CVaR \cite{rockafellar2000} y máximo Sharpe regularizado. El
Expected Shortfall discreto usa masa fraccionaria en el cuantil; así se evita sesgo
cuando una probabilidad cruza el nivel $\alpha$.

El backtest deja derivar los pesos, registra turnover y descuenta costos en puntos
base. MM y los controles generativos se evalúan como asignaciones estáticas ajustadas
una sola vez con la muestra inicial; Equal Weight, inverse variance y HRP incluyen
además versiones trimestrales con ventana expansiva que usan solo observaciones
anteriores al rebalanceo. La evidencia estadística se expresa mediante intervalos
moving-block bootstrap del diferencial de Sharpe frente a Equal Weight.

\section{Resultados de la corrida 2026-08-10}

La corrida vigente contiene 993 observaciones rolling de cinco días para estimación
y 121 observaciones no solapadas en validación. Se corrigió una tolerancia L-BFGS
incompatible con la escala de $F$: SciPy comparaba $10^{-9}$ contra una escala mínima
unitaria pese a que el objetivo era del orden de $10^{-7}$. Con el parámetro
$x_{\mathrm{ftol}}=10^{-12}$, los cinco starts alcanzan el criterio de convergencia
en 5--6 iteraciones. El start publicado es el menor $G$ entre los convergentes que
además cumplen estacionariedad: obtiene $F=1.06\times10^{-10}$,
$G=6.34\times10^{-6}$ y $N_{\mathrm{eff}}=166.2$ de 500 escenarios. Los residuos
infinito son $1.12\times10^{-5}$ para $X$ y $8.02\times10^{-5}$ para el gradiente
tangente de $p$.

\begin{table}[ht]
\centering
\small
\begin{tabular}{lrrrr}
\toprule
Modelo & CRPS & Energy & Variogram & Excepciones 5\% \\
\midrule
Student-$t$ & 0.020633 & 0.100045 & 0.611200 & 2 \\
Gaussiano & 0.020668 & 0.099989 & 0.623808 & 0 \\
MM & 0.020738 & 0.100094 & 0.618069 & 1 \\
Histórico EWMA & 0.020734 & 0.100132 & 0.617487 & 3 \\
\bottomrule
\end{tabular}
\caption{Validación probabilística; menor score es mejor. ``Excepciones'' cuenta
activos cuyo test de cobertura VaR rechaza al 5\%.}
\end{table}

MM no lidera los scores: Student-$t$ obtiene el mejor CRPS y Gaussiano el mejor
Energy Score. En la
evaluación económica estática, inverse variance alcanza Sharpe 2.032 y MM de mínima
varianza 1.959, ambos netos de 20 bps por unidad de turnover. El intervalo bootstrap
95\% del diferencial de Sharpe anualizado de MM-MinVariance frente a Equal Weight
walk-forward es $[-0.152,\,0.334]$; contiene cero. Por tanto, la evidencia actual no
respalda superioridad estadística ni económica de MM. Es un resultado de desarrollo,
no una conclusión confirmatoria.

\subsection*{Incertidumbre de las diferencias}

La diferencia se define como MM menos control, por lo que un valor positivo es
desfavorable a MM. El contraste primario CRPS contra Gaussiano cruza cero y ninguno
de los nueve contrastes del universo completo sobrevive Holm.

\begin{table}[ht]
\centering
\small
\begin{tabular}{llrrr}
\toprule
Universo & Contraste & Diferencia & IC95 & p Holm \\
\midrule
Completo & CRPS vs Gaussiano & +0.000070 & [-0.000044,0.000162] & 1.000 \\
Completo & CRPS vs Student-t & +0.000105 & [0.000007,0.000180] & 0.262 \\
Completo & Variogram vs Student-t & +0.006868 & [0.000864,0.012811] & 0.198 \\
Líquido & CRPS vs Gaussiano & +0.000070 & [-0.000015,0.000137] & 0.693 \\
Líquido & Variogram vs hist. & +0.001660 & [0.000577,0.002608] & 0.0108 \\
\bottomrule
\end{tabular}
\caption{Bootstrap temporal pareado. Holm se aplica a nueve contrastes dentro de
cada universo; la tabla muestra el contraste primario y resultados representativos.}
\end{table}

La única diferencia que sobrevive Holm aparece en la sensibilidad líquida y es
contraria a MM: su Variogram Score es peor que el histórico. El hallazgo sigue siendo
de desarrollo porque el periodo ya fue observado.

\section{Robustez temporal rolling-origin}

Los cuatro folds contienen 169 ventanas OOS: 49 en 2023, 49 en 2024, 49 en 2025 y
22 en 2026-H1. Los cuatro modelos se recalibran en cada origen con idénticos
hiperparámetros. Todos los starts publicados cumplen convergencia y estacionariedad.

\begin{table}[ht]
\centering
\small
\begin{tabular}{lrrrr}
\toprule
Modelo & CRPS pooled & Energy pooled & Variogram pooled & Victorias CRPS \\
\midrule
Student-t & 0.020858 & 0.100672 & 0.642320 & 2/4 \\
Gaussiano & 0.020935 & 0.100850 & 0.663707 & 1/4 \\
Histórico EWMA & 0.020969 & 0.101302 & 0.654289 & 1/4 \\
MM & 0.020972 & 0.101139 & 0.657408 & 0/4 \\
\bottomrule
\end{tabular}
\caption{Evaluación rolling-origin expansiva; menor score es mejor.}
\end{table}

El contraste primario pooled CRPS MM menos Gaussiano se estima en +0.000036.
Su IC95 va de -0.000041 a 0.000118 y el p Holm es 0.682. Student-$t$ supera a MM
en CRPS: diferencia +0.000114, IC95 de 0.000058 a 0.000183 y p Holm de 0.004.
MM sí mejora al histórico en Energy Score y al Gaussiano en Variogram Score. Por
tanto, la conclusión depende de la propiedad distributiva y no existe superioridad
general de MM.

\section{Robustez frente a iliquidez}

Como sensibilidad separada, se congeló antes de recalibrar una regla que conserva
activos con tasa de retornos cero IS no mayor a 3\% y racha IS no mayor a cinco.
La regla excluye AGUAS-A, ANDINA-B, ITAUCL y PARAUCO y deja 11 activos. Ninguna
métrica OOS interviene en la selección; una prueba automatizada muta esas columnas y
exige que la decisión permanezca invariante. El universo principal de 15 activos no
se reemplaza.

\begin{table}[ht]
\centering
\small
\begin{tabular}{lrrrr}
\toprule
Modelo & CRPS líquido & Energy líquido & Rango completo & Rango líquido \\
\midrule
Student-$t$ & 0.021250 & 0.087987 & 1 & 1 \\
Gaussiano & 0.021290 & 0.088061 & 2 & 2 \\
MM & 0.021360 & 0.088244 & 4 & 3 \\
Histórico EWMA & 0.021365 & 0.088281 & 3 & 4 \\
\bottomrule
\end{tabular}
\caption{Sensibilidad en 11 activos. Los niveles entre universos de distinta
dimensión no son directamente comparables; el rango es interno a cada universo.}
\end{table}

MM supera al histórico líquido por solo $4\times10^{-6}$ de CRPS, sin evidencia
estadística, mientras Student-$t$ y Gaussiano conservan los dos primeros lugares.
MM-MinVariance obtiene Sharpe 1.748
y drawdown $-12.98\%$, frente a 1.959 y $-13.80\%$ con 15 activos. Su intervalo
bootstrap 95\% contra Equal Weight walk-forward es $[-0.165,\,0.428]$ e incluye cero.
La reducción del drawdown coincide con una caída general de Sharpe: el filtro no
demuestra una mejora económica ni cambia la conclusión principal.

\section{Reglas de interpretación}

MM solo se declarará superior si mejora scores propios fuera de muestra, mantiene
cobertura de cola razonable y produce una mejora económica neta cuyo intervalo de
incertidumbre sea compatible con la afirmación. Un resultado negativo también es un
resultado de investigación válido.

Los artefactos previos a esta versión son legacy: cambió el gradiente, la escala de
momentos, el término de covarianza, el criterio de convergencia, el EWMA y los
benchmarks. Por ello sus valores numéricos no se mezclan con una nueva corrida.

\section{Reproducibilidad}

Cada corrida final contiene configuración, versiones, hash de código y SHA-256 de
artefactos. La suite automatizada verifica gradientes, monotonía de $G$, composición
de retornos, parametrización de benchmarks, scores, restricciones de portafolio y
ausencia de look-ahead en el motor walk-forward y en el filtro de liquidez, agregación
de pérdidas pareadas, folds rolling-origin, ajuste Holm y selección de starts
estacionarios. Los 45 tests
verdes prueban contratos de
software; no prueba rentabilidad futura.

Además, cada frontera del pipeline mantiene un manifiesto de linaje con hashes de
entradas, código relevante y salidas. Transformación, MM, benchmarks, evaluación,
portafolios y backtest se niegan a consumir un upstream cuyo contenido cambió. El
snapshot final solo se crea si los nueve manifiestos, incluidas la descarga, la
sensibilidad de liquidez y la robustez rolling-origin, permanecen vigentes; así se evita
combinar silenciosamente escenarios nuevos con scores o decisiones anteriores.

\Needspace{12\baselineskip}
\small
\begin{thebibliography}{99}
\bibitem{hoyland2001} Høyland, K. y Wallace, S. W. (2001). Generating scenario trees for multistage decision problems. \textit{Management Science}, 47(2), 295--307.
\bibitem{ponomareva2015} Ponomareva, K., Roman, D. y Date, P. (2015). An algorithm for moment-matching scenario generation with application to financial portfolio optimisation. \textit{European Journal of Operational Research}, 240(3), 678--687.
\bibitem{contreras2018} Contreras, J. P., Bosch, P. y Herrera, M. (2018). Comment on ``An algorithm for moment-matching scenario generation with application to financial portfolio optimization''. \textit{European Journal of Operational Research}, 269(3), 1180--1184.
\bibitem{mehrotra2013} Mehrotra, S. y Papp, D. (2013). Generating moment matching scenarios using optimization techniques. \textit{SIAM Journal on Optimization}, 23(2).
\bibitem{gneiting2007} Gneiting, T. y Raftery, A. E. (2007). Strictly proper scoring rules, prediction, and estimation. \textit{JASA}, 102(477), 359--378.
\bibitem{kupiec1995} Kupiec, P. H. (1995). Techniques for verifying the accuracy of risk measurement models. \textit{Journal of Derivatives}, 3(2), 73--84.
\bibitem{christoffersen1998} Christoffersen, P. F. (1998). Evaluating interval forecasts. \textit{International Economic Review}, 39(4), 841--862.
\bibitem{ledoitwolf2004} Ledoit, O. y Wolf, M. (2004). A well-conditioned estimator for large-dimensional covariance matrices. \textit{Journal of Multivariate Analysis}, 88(2), 365--411.
\bibitem{lopez2016} López de Prado, M. (2016). Building diversified portfolios that outperform out-of-sample. \textit{Journal of Portfolio Management}, 42(4), 59--69.
\bibitem{rockafellar2000} Rockafellar, R. T. y Uryasev, S. (2000). Optimization of conditional value-at-risk. \textit{Journal of Risk}, 2(3), 21--41.
\bibitem{lesmond1999} Lesmond, D. A., Ogden, J. P. y Trzcinka, C. A. (1999). A new estimate of transaction costs. \textit{Review of Financial Studies}, 12(5), 1113--1141.
\bibitem{diebold1995} Diebold, F. X. y Mariano, R. S. (1995). Comparing predictive accuracy. \textit{Journal of Business \& Economic Statistics}, 13(3), 253--263.
\bibitem{kunsch1989} Künsch, H. R. (1989). The jackknife and the bootstrap for general stationary observations. \textit{The Annals of Statistics}, 17(3), 1217--1241.
\bibitem{holm1979} Holm, S. (1979). A simple sequentially rejective multiple test procedure. \textit{Scandinavian Journal of Statistics}, 6(2), 65--70.
\bibitem{tashman2000} Tashman, L. J. (2000). Out-of-sample tests of forecasting accuracy: an analysis and review. \textit{International Journal of Forecasting}, 16(4), 437--450.
\bibitem{hewamalage2023} Hewamalage, H., Ackermann, K. y Bergmeir, C. (2023). Forecast evaluation for data scientists: common pitfalls and best practices. \textit{Data Mining and Knowledge Discovery}, 37, 788--832.
\end{thebibliography}
\normalsize

\end{document}
